\section{The induced sliding block code}
\label{sec:sequence}

For each $m\ge1$, the window fold induces a sliding block code on the full two-sided shift:
\begin{equation}\label{eq:Phi}
(\Phi_m(\eta))_t
:=
\Fold_m(\eta_{t-m+1}\cdots \eta_t),
\qquad \eta\in\{0,1\}^{\ZZ},\ t\in\ZZ.
\end{equation}
Its image shift is
\[
Y_m:=\Phi_m(\{0,1\}^{\ZZ})\subset X_m^{\ZZ},
\]
equipped with the left shift on the $X_m$-alphabet. The natural question is whether the microstate sequence can be recovered locally from the sequence of folded windows.

\subsection{A graph presentation}

\begin{proposition}\label{prop:sofic}
For each $m\ge1$, the image shift $Y_m$ is sofic and admits an explicit labeled-graph presentation.
\end{proposition}

\begin{proof}
Let $V_m:=\{0,1\}^{m-1}$. For $v\in V_m$ and $b\in\{0,1\}$, define the successor vertex $v'$ by shifting and appending $b$,
\[
v'=\text{suffix}_{m-1}(vb),
\]
and label the directed edge $v\to v'$ by
\[
\lambda_m(vb):=\Fold_m(vb)\in X_m.
\]
Denote the resulting labeled graph by $G_m$. Let $\Sigma_m$ be its edge shift and $\Lambda_m:\Sigma_m\to X_m^{\ZZ}$ the label map. By construction $\Lambda_m(\Sigma_m)=Y_m$, so $Y_m$ is sofic.
\end{proof}

\subsection{The sharp ambiguity bound}

For a block $w=(w_1,\ldots,w_L)\in\{0,1\}^L$, $L\ge m$, write
\[
\mathcal{W}_{m,L}(w)
:=
\bigl(\Fold_m(w_i\cdots w_{i+m-1})\bigr)_{i=1}^{L-m+1}
\]
for its sequence of folded length-$m$ windows.

\begin{theorem}[Sharp block separation]\label{thm:block-separation}
Let $m\ge3$. If $L\ge2m-1$, then the map
\[
\mathcal{W}_{m,L}:\{0,1\}^L\longrightarrow X_m^{L-m+1}
\]
is injective. The threshold is sharp: for every $m\ge3$, distinct blocks of length $2m-2$ have the same folded-window sequence.
\end{theorem}

\begin{proof}
Set $M:=\Fib_{m+2}$. A length-$m$ word has value at most
\[
\sum_{k=1}^m\Fib_{k+1}=\Fib_{m+3}-2
=M+\Fib_{m+1}-2.
\]
Its normalized expansion therefore has at most one digit beyond the window, at position $m+1$. When that digit occurs, the remainder is $N-M<\Fib_{m+1}$, so it cannot use position $m$; the retained digits are consequently the Zeckendorf expansion of $N\bmod M$. Thus two length-$m$ words have the same fold precisely when their values are congruent modulo $M$.

It is enough to prove injectivity for $L=2m-1$: if two longer blocks differ, some length-$2m-1$ subblock containing a differing coordinate also differs. Let $u,v\in\{0,1\}^{2m-1}$ have the same window sequence, put $d_j:=u_j-v_j\in\{-1,0,1\}$, and define
\[
S_i:=\sum_{k=1}^m d_{i+k-1}\Fib_{k+1},
\qquad 1\le i\le m.
\]
The collision criterion gives $S_i\equiv0\pmod M$. For $1\le i\le m-2$, the Fibonacci recurrence cancels all interior coefficients in three consecutive sums:
\begin{equation}\label{eq:window-boundary-identity}
S_i-S_{i+1}-S_{i+2}
=d_i+d_{i+1}-M d_{i+m}-\Fib_{m+1}d_{i+m+1}.
\end{equation}

Suppose first that $m\ge4$. Reducing \eqref{eq:window-boundary-identity} modulo $M$ gives
\[
d_i+d_{i+1}-\Fib_{m+1}d_{i+m+1}\equiv0\pmod M.
\]
Its absolute value is at most $2+\Fib_{m+1}<M$, because $\Fib_m>2$. It is therefore zero. Since $\Fib_{m+1}>2$, this forces
\begin{equation}\label{eq:window-boundary-consequence}
d_{i+m+1}=0,
\qquad d_{i+1}=-d_i
\qquad(1\le i\le m-2).
\end{equation}
Hence $d_{m+2}=\cdots=d_{2m-1}=0$, while $d_1,\ldots,d_{m-1}$ alternate. The last window now satisfies
\[
S_m=d_m+2d_{m+1}\equiv0\pmod M.
\]
Since $|S_m|\le3<M$, it vanishes; the digit bounds then give $d_m=d_{m+1}=0$. Finally,
\[
S_1=d_1\sum_{j=1}^{m-1}(-1)^{j-1}\Fib_{j+1}
=d_1(-1)^{m-2}\Fib_{m-1}.
\]
This is divisible by $M$ and has absolute value smaller than $M$, so $d_1=0$. Equation \eqref{eq:window-boundary-consequence} now gives $d=0$, and $u=v$.

For $m=3$, one has $M=5$, and reduction of \eqref{eq:window-boundary-identity} gives
\[
d_1+d_2-3d_5\equiv0\pmod5.
\]
If $d_5=0$, then $d_2=-d_1$; the last window gives $d_3+2d_4=0$, hence $d_3=d_4=0$, and the first gives $-d_1=0$. If $d_5=1$, the congruence forces $d_1=d_2=-1$. The first window then forces $d_3=1$, whereas the last window has value $4+2d_4\in\{2,4,6\}$, none divisible by $5$. The case $d_5=-1$ follows by negation. Thus $d=0$ also when $m=3$.

For sharpness, take length $2m-2$ blocks $u\ne v$ defined by
\[
u_{m+1}=1,
\qquad
v_{m-1}=v_m=1,
\]
with every other digit zero. The first window-value difference is
$-\Fib_m-\Fib_{m+1}=-M$. Every later window contains all three nonzero coordinates, and its value difference is zero by the Fibonacci recurrence. Hence $\mathcal{W}_{m,2m-2}(u)=\mathcal{W}_{m,2m-2}(v)$.
\end{proof}

\subsection{Local decoding}

Theorem~\ref{thm:block-separation} says that $m$ consecutive labels determine their entire length-$(2m-1)$ lift. Recovering only the current digit requires fewer labels.

\begin{theorem}[Exact causal decoder]\label{thm:window-decoder}
Fix $m\ge3$. Any two lifts of the same three consecutive symbols in $Y_m$ have the same terminal digit. Consequently there is a well-defined decoder
\[
\psi_m:\mathcal{B}_3(Y_m)\to\{0,1\},
\qquad
\eta_t=\psi_m(y_{t-2},y_{t-1},y_t)
\]
whenever $y=\Phi_m(\eta)$. No decoder using only $(y_{t-1},y_t)$ can recover $\eta_t$. Thus the exact zero-anticipation memory of the inverse is $2$, independently of $m$.
\end{theorem}

\begin{proof}
Take two length-$(m+2)$ lifts of the same three labels and write their digitwise difference as $d_1,\ldots,d_{m+2}\in\{-1,0,1\}$. The three window differences $S_1,S_2,S_3$ from the proof of \Cref{thm:block-separation} are divisible by $M=\Fib_{m+2}$. Applying \eqref{eq:window-boundary-identity} at $i=1$ and reducing modulo $M$ gives
\[
d_1+d_2-\Fib_{m+1}d_{m+2}\equiv0\pmod M.
\]
If $m\ge4$, the absolute value of the left-hand side is at most $2+\Fib_{m+1}<M$, so it vanishes. Since $\Fib_{m+1}>2$, this forces $d_{m+2}=0$. For $m=3$, three labels correspond to a length-$5=2m-1$ lift, which is unique by \Cref{thm:block-separation}. Thus the terminal digit is determined in every case.

To prove optimality, take length-$(m+1)$ blocks $u\ne v$ with
\[
u_{m+1}=1,
\qquad
v_{m-1}=v_m=1,
\]
and all other digits zero. In the first window the values differ by
$\Fib_m+\Fib_{m+1}=M$, so their folds agree. In the second window the two values are equal by the same Fibonacci recurrence. Hence the two labels agree, while the terminal digits are $u_{m+1}=1$ and $v_{m+1}=0$. Two labels therefore do not suffice.
\end{proof}

\begin{remark}\label{rem:window-bound-essential}
The two sharp thresholds concern different recovery problems. A sequence of $m$ labels is necessary to determine the whole length-$(2m-1)$ lift, whereas three labels determine its terminal digit. Thus sharp whole-block reconstruction at length $2m-1$ does not determine the inverse memory.
\end{remark}

\subsection{The threshold \texorpdfstring{$m\ge3$}{m >= 3}}

\begin{theorem}\label{thm:conjugacy}
For $m\ge3$, the sliding block code $\Phi_m:\{0,1\}^{\ZZ}\to Y_m$ is a topological conjugacy: there exists a sliding block inverse $\Psi_m:Y_m\to \{0,1\}^{\ZZ}$ such that
\[
\Psi_m\circ \Phi_m=\mathrm{id}_{\{0,1\}^{\ZZ}},
\qquad
\Phi_m\circ \Psi_m=\mathrm{id}_{Y_m}.
\]
\end{theorem}

\begin{proof}
For $y\in Y_m$, define
\[
(\Psi_m(y))_t
:=\psi_m(y_{t-2},y_{t-1},y_t),
\]
where $\psi_m$ is the decoder from \Cref{thm:window-decoder}. This is a sliding block code and $\Psi_m\circ\Phi_m=\mathrm{id}$. If $y\in Y_m$, write $y=\Phi_m(\eta)$; then $\Psi_m(y)=\eta$, so $\Phi_m\circ\Psi_m(y)=y$. Hence $\Psi_m$ is the desired inverse.
\end{proof}

\begin{proposition}\label{thm:m2}
When $m=2$, the only failure of injectivity is that the two fixed microstate sequences $0^\ZZ$ and $1^\ZZ$ have the same image $(00)^\ZZ$. Away from this pair, $\Phi_2$ is injective.
\end{proposition}

\begin{proof}
Appendix~\ref{app:sofic} identifies the only off-diagonal bi-infinite cycle in the fiber product for $m=2$: it is the constant label sequence $(00)^\ZZ$, lifted once by $0^\ZZ$ and once by $1^\ZZ$. Every other label sequence contains $01$ or $10$, and those labels determine the adjacent microstate bits uniquely. Hence $(00)^\ZZ$ is the unique non-injective fiber.
\end{proof}
